\documentclass[14pt,a4paper]{extarticle}

\usepackage[T2A]{fontenc}
\usepackage[utf8]{inputenc}
\usepackage[russian,english]{babel}

\usepackage{geometry}
\geometry{left=30mm,right=15mm,top=20mm,bottom=20mm}

\usepackage{setspace}
\onehalfspacing

\usepackage{indentfirst}
\setlength{\parindent}{1.25cm}

\usepackage{enumitem}
\setlist[itemize]{leftmargin=1.25cm}
\setlist[enumerate]{leftmargin=1.25cm}

\usepackage{longtable}
\usepackage{ltablex}
\usepackage{array}
\usepackage{tabularx}
\usepackage{booktabs}
\usepackage{multirow}
\usepackage{ragged2e}
\usepackage{xurl}
\usepackage{hyperref}
\hypersetup{hidelinks}
\keepXColumns

\setlength{\emergencystretch}{3em}
\setlength{\tabcolsep}{3pt}
\renewcommand{\arraystretch}{1.2}
\newcolumntype{L}[1]{>{\RaggedRight\arraybackslash}p{#1}}
\newcolumntype{C}[1]{>{\Centering\arraybackslash}p{#1}}
\newcolumntype{Y}{>{\RaggedRight\arraybackslash}X}
\sloppy

\begin{document}

\begin{titlepage}
\begin{center}
Министерство науки и высшего образования Российской Федерации\\
Национальный исследовательский ядерный университет «МИФИ»\\
Институт интеллектуальных кибернетических систем\\
Направление подготовки «Программная инженерия»
\end{center}

\vspace{2cm}

\begin{center}
\textbf{\Large ПОЯСНИТЕЛЬНАЯ ЗАПИСКА}\\
\vspace{0.5cm}
\textbf{\large к проекту «Агентный аналитик данных»}
\end{center}

\vspace{2cm}

\noindent
\textbf{Команда проекта:}\\
Холодков Антон Евгеньевич\\
Голод Герман Матвеевич\\
Малышев Артём Вячеславович\\
Песковацков Илья Александрович

\vspace{1cm}

\noindent
\textbf{Учебная группа:} Б24-517\\
\textbf{Курс:} 2\\
\textbf{Вид работы:} проектная практика

\vfill

\begin{center}
Москва --- 2026
\end{center}
\end{titlepage}

\renewcommand{\contentsname}{Содержание}
\tableofcontents
\newpage

\section{Введение}

Настоящая пояснительная записка подготовлена для проекта «Агентный аналитик данных», выполняемого командой студентов 2 курса группы Б24-517 направления «Программная инженерия искусственного интеллекта» НИЯУ «МИФИ». Документ ориентирован на структуру и стиль изложения, принятые для программной документации по ГОСТ 19.404--79, и фиксирует назначение проекта, состав разрабатываемой системы, ключевые технические решения, организацию работ, сроки выполнения и ожидаемые результаты.

Проект направлен на создание системы агентной аналитики данных, способной принимать пользовательскую задачу, подключаться к источникам данных, выполнять маршрутизацию между инструментами и подагентами, проводить базовую обработку наборов данных, формировать аналитический результат и обеспечивать обезличивание чувствительных данных при работе с реальными источниками.

Принцип разработки системы заключается в том, чтобы не использовать агентный подход там, где требуемый результат может быть получен обычной программной логикой предсказуемо, воспроизводимо и с меньшими издержками. Агенты и LLM-компоненты применяются только в тех частях решения, где они действительно дают выигрыш по гибкости обработки, интерпретации запроса или формированию аналитического вывода.

\section{Основание для разработки}

Основанием для выполнения работы является задание в рамках проектной практики по командному проекту «Агентный аналитик данных». В переданных требованиях зафиксированы следующие ключевые положения:
\begin{itemize}
    \item целевая архитектура строится вокруг единой точки входа с инструментами и агентами;
    \item в качестве базового состояния системы рассматривается оркестратор;
    \item в состав решения входят работа с датасетами, инструменты аналитики и модуль Data Anonymizer;
\end{itemize}

\section{Назначение и цели проекта}

\subsection{Назначение системы}

Разрабатываемая система предназначена для выполнения аналитических задач по пользовательскому запросу с использованием набора специализированных инструментов и подагентов под управлением центрального оркестратора. Система должна обеспечить:
\begin{itemize}
    \item прием и интерпретацию пользовательского запроса;
    \item подключение к источникам данных;
    \item предварительный просмотр и анализ структуры датасета;
    \item проверку качества данных;
    \item преобразование данных и выполнение аналитических операций;
    \item формирование итогового ответа в пригодном для пользователя виде;
    \item маскирование чувствительных данных перед передачей информации внешним LLM-компонентам.
\end{itemize}

\subsection{Цели разработки}

Целями разработки являются:
\begin{enumerate}
    \item создание работоспособной системы «Агентный аналитик данных»;
    \item реализация оркестратора как центрального узла управления;
    \item реализация набора инструментов и подагентов для работы с датасетами;
    \item включение модуля Data Anonymizer для безопасной работы с чувствительными данными;
    \item получение экспериментальных результатов по качеству работы и производительности системы;
    \item подготовка оформленных материалов по итогам замеров и анализа результатов.
\end{enumerate}

\section{Краткая характеристика объекта разработки}

Проект «Агентный аналитик данных» представляет собой программную систему агентного типа, в которой один верхнеуровневый агент выполняет координацию специализированных исполнительных компонентов. Система ориентирована на задачи прикладной аналитики: чтение данных из таблиц и файлов, базовую очистку, преобразование, построение сводных результатов и формирование аналитического ответа.

Выбранная архитектура подразумевает максимально возможный детерминизм принятия решений и использует вероятностные методы(агенты) там, где это оправдано.

\section{Требования к разрабатываемой системе}

\subsection{Функциональные требования}

Система должна обеспечивать следующие функции:
\begin{enumerate}
    \item прием пользовательского запроса на естественном языке;
    \item выбор допустимого сценария обработки запроса;
    \item доступ к источникам данных в файлах и табличных форматах;
    \item получение описания структуры набора данных;
    \item проверку качества данных;
    \item выполнение преобразований данных;
    \item вызов аналитических инструментов;
    \item построение текстового результата и, при необходимости, структурированного артефакта;
    \item обезличивание чувствительных данных до передачи во внешние компоненты.
\end{enumerate}

\subsection{Нефункциональные требования}

К системе предъявляются следующие нефункциональные требования:
\begin{itemize}
    \item идемпотетность аналитических операций;
    \item воспроизводимость сценариев выполнения;
    \item модульность архитектуры;
    \item возможность расширения состава инструментов без полной переработки ядра;
    \item ограничение небезопасных операций;
    \item пригодность для демонстрации на учебных примерах и датасетах;
    \item наличие измеримых характеристик по времени выполнения и качеству результата.
\end{itemize}

\section{Обоснование выбранных технических решений}

\subsection{Выбор оркестратора в качестве ядра системы}

Базовым архитектурным решением является использование оркестратора. Это решение выбрано по следующим причинам:
\begin{itemize}
    \item соответствует переданным требованиям, где именно оркестратор отмечен как основа;
    \item озволяет централизованно управлять ходом выполнения задачи;
    \item упрощает тестирование, так как маршрут обработки запроса можно фиксировать и воспроизводить;
    \item позволяет отделить этап построения надежного ядра от последующего этапа интеллектуальной автоматизации;
    \item уменьшает риск циклических вызовов, лишних обращений к инструментам и труднообъяснимых результатов.
\end{itemize}

Текущая реализация предполагает базовый набор пользовательских сценариев и стремится им соответстовать.

\subsection{Принцип минимального использования агентов}

При разработке системы принимается принцип: агентный подход должен использоваться только там, где он действительно необходим. Если результат может быть получен с помощью обычного алгоритма, жестко заданного пайплайна, правил валидации, SQL-запроса, преобразования таблицы или иного программного механизма с предсказуемым поведением, то предпочтение отдается именно такому способу реализации.

Данный принцип выбран по следующим причинам:
\begin{itemize}
    \item снижает сложность системы;
    \item повышает воспроизводимость результатов;
    \item уменьшает вычислительные затраты;
    \item позволяет применять агентные компоненты только в задачах, где нужна гибкая интерпретация или генерация.
\end{itemize}

\subsection{Выбор агентно-инструментальной схемы}

В архитектуре сочетаются два уровня исполнения:
\begin{enumerate}
    \item прямые инструменты, доступные верхнему агенту;
    \item специализированные агентные модули, вызываемые как инструменты более высокого уровня.
\end{enumerate}

Такой подход позволяет соблюсти баланс между простотой и расширяемостью. Простые и безопасные операции, например просмотр схемы, загрузка артефакта или предварительный просмотр данных, доступны напрямую. Более сложные и доменно-специфические задачи инкапсулируются в отдельных модулях. В результате ядро остается компактным, а логика работы отдельных подсистем становится локализованной.

\subsection{Выбор Data Anonymizer как опциональной части системы}

Модуль обезличивания включается в состав системы, поскольку работа с реальными данными без предварительного маскирования чувствительной информации ограничивает возможность практического применения решения. Даже если большая часть демонстраций будет проводиться на учебных наборах, наличие Data Anonymizer делает архитектуру завершенной и показывает готовность проекта к переходу от демонстрационного режима к реальным данным.

На этапе реализации модуль обезличивания должен решать минимально необходимую задачу:
\begin{itemize}
    \item обнаружение типовых чувствительных сущностей;
    \item замена найденных значений на безопасные токены;
    \item сохранение таблицы соответствия для обратного восстановления при необходимости;
    \item предотвращение передачи открытых чувствительных данных во внешние LLM-компоненты.
\end{itemize}

\subsection{Минимально реализуемый вариант генерации инструментов и промптов}

В проект закладывается минимально реализуемый вариант:
\begin{itemize}
    \item полуавтоматическое подключение новых инструментов через декларативное описание;
    \item реестр инструментов с полями: имя, назначение, допустимые входы, формат выхода, ограничения вызова;
    \item шаблонный компоновщик промптов, который подставляет в заранее подготовленные шаблоны контекст задачи, описание датасета и тип требуемого результата;
    \item небольшой набор фиксированных промпт-режимов, например: анализ качества данных, сводная аналитика, преобразование таблицы, формирование отчета;
    \item ручное подтверждение или ручная регистрация нового инструмента разработчиком.
\end{itemize}

\section{Описание архитектуры системы}

\subsection{Общая схема}

Логика работы системы строится по следующей схеме:
\begin{center}
\parbox{0.95\textwidth}{\centering Пользователь $\rightarrow$ Supervisor / Chief Analyst Agent $\rightarrow$ инструменты и подагенты $\rightarrow$ формирование результата}
\end{center}

Верхний агент интерпретирует запрос, определяет сценарий обработки, выбирает необходимые компоненты, контролирует промежуточные артефакты и формирует итоговый ответ.

\subsection{Роль верхнего агента}

Supervisor / Chief Analyst Agent является центральным узлом управления. Он выполняет следующие функции:
\begin{itemize}
    \item интерпретация пользовательского запроса;
    \item декомпозиция задачи на этапы;
    \item выбор между прямым вызовом инструмента и вызовом агентного модуля;
    \item контроль последовательности выполнения;
    \item консолидация промежуточных результатов;
    \item подготовка итогового ответа.
\end{itemize}

\subsection{Прямые инструменты верхнего агента}

Наиболее безопасные и универсальные операции предоставляются напрямую:
\begin{itemize}
    \item реестр инструментов/MCP/A2A;
    \item просмотр метаданных;
    \item загрузка и сохранение артефактов;
    \item предварительный просмотр датасета;
    \item инспекция схемы;
    \item разбор и валидация JSON;
    \item предварительная работа с Excel;
    \item безопасный предпросмотр SQL;
    \item форматирование результата.
\end{itemize}

\subsection{Специализированные модули системы}

В составе системы целесообразно выделить ограниченное число подагентов, а остальную функциональность оформить как инструменты или MCP-компоненты.

\subsubsection{Подагенты системы}

В качестве отдельных подагентов целесообразно реализовать следующие компоненты:

\begin{enumerate}
    \item \textbf{Data Understanding Agent} --- анализирует запрос пользователя, соотносит его со структурой доступных данных и определяет общий сценарий обработки.

    \item \textbf{Analytical Reasoning Agent} --- выбирает логику аналитической обработки: какие шаги, инструменты и операции нужны для решения задачи.

    \item \textbf{Reporting Agent} --- собирает промежуточные результаты и формирует итоговый ответ или структуру отчета.
\end{enumerate}

\subsubsection{Инструменты и MCP-компоненты}

Следующие функции целесообразно оформить как инструменты или MCP-компоненты:

\begin{enumerate}
    \item \textbf{Data Source MCP / Tool} --- загрузка файлов, таблиц и иных источников данных.
    \item \textbf{Schema Inspection Tool} --- просмотр структуры датасета, типов полей и примеров значений.
    \item \textbf{Data Quality Tool} --- поиск пропусков, дубликатов и других проблем качества данных.
    \item \textbf{Transformation Tool} --- фильтрация, агрегация, объединение, сортировка и другие преобразования.
    \item \textbf{SQL Execution MCP / Tool} --- безопасное выполнение аналитических SQL-запросов.
    \item \textbf{Python Analytics Tool} --- вычисления и аналитическая обработка данных.
    \item \textbf{Visualization Tool} --- построение графиков, диаграмм и иных визуальных представлений.
    \item \textbf{Data Anonymizer MCP / Tool} --- обезличивание чувствительных данных и обратное восстановление при необходимости.
    \item \textbf{Artifact Storage Tool} --- сохранение промежуточных и итоговых артефактов.
\end{enumerate}

\subsubsection{Итоговый принцип разбиения}

Таким образом, подагенты должны отвечать за интерпретацию задачи, выбор логики обработки и формирование ответа, а операции чтения данных, проверки качества, преобразования, вычислений, SQL-аналитики, визуализации и обезличивания следует реализовать как инструменты или MCP-компоненты.

\subsection{Принцип прохождения данных}

Общий поток обработки данных предлагается организовать следующим образом:
\begin{enumerate}
    \item пользователь задает аналитическую задачу;
    \item оркестратор определяет нужный сценарий;
    \item данные загружаются из источника;
    \item выполняется базовый просмотр и анализ качества;
    \item при необходимости выполняется обезличивание;
    \item данные передаются в модуль преобразования и аналитики;
    \item полученные результаты сохраняются как артефакты;
    \item формируется итоговый ответ.
\end{enumerate}

\section{План реализации проекта}

\subsection{Общий график}

Срок выполнения работ делится на два крупных этапа:
\begin{itemize}
    \item 5 недель --- разработка;
    \item 3 недели --- замеры, анализ и оформление результатов.
\end{itemize}

\subsection{План работ по неделям}

\begin{tabularx}{\textwidth}{|C{1.8cm}|L{4.8cm}|Y|}
\hline
\textbf{Неделя} & \textbf{Этап} & \textbf{Содержание работ} \\
\hline
1 & Проектирование ядра & Уточнение архитектуры системы, фиксация сценариев работы оркестратора, описание интерфейсов модулей, определение форматов артефактов и контрактов между компонентами. \\
\hline
2 & Реализация базового контура & Разработка каркаса Supervisor / Chief Analyst Agent, реестра инструментов, базовых direct tools, механизма вызова агентных модулей. \\
\hline
3 & Работа с датасетами & Реализация Data Source Agent, Data Quality Agent и Transformation Agent, подключение типовых форматов данных, тестирование на примерах. \\
\hline
4 & Аналитика и отчетность & Реализация SQL Analysis Agent, Python Analytics Agent, Reporting Agent, настройка прохождения артефактов, стабилизация основных пользовательских сценариев. \\
\hline
5 & Обезличивание и интеграция & Реализация Data Anonymizer, включение маскирования в общий конвейер, интеграционное тестирование системы, исправление критических дефектов. \\
\hline
6 & Замеры производительности & Проведение серий замеров на подготовленных датасетах, сбор метрик по времени выполнения, устойчивости сценариев, полноте результатов и корректности маскирования. \\
\hline
7 & Анализ результатов & Сравнение сценариев, интерпретация результатов, выявление узких мест, формирование выводов о сильных и слабых сторонах системы. \\
\hline
8 & Оформление материалов & Подготовка итоговой документации, таблиц, диаграмм, выводов, формирование презентационных и отчетных материалов. \\
\hline
\end{tabularx}

\section{Распределение работ в команде}

С учетом переданного контекста и необходимости закрепить зоны ответственности предлагается следующее распределение.

\begin{tabularx}{\textwidth}{|L{4.6cm}|Y|}
\hline
\textbf{Участник} & \textbf{Предполагаемая зона ответственности} \\
\hline
Песковацков Илья Александрович & Архитектурная координация, проектирование и реализация оркестратора, формализация взаимодействия модулей, описание сценариев выполнения, участие в интеграции и итоговом оформлении результатов. \\
\hline
Малышев Артём Вячеславович & Реализация компонентов визуализации, интерпретации данных и формирования итогового представления результатов для пользователя. \\
\hline
Холодков Антон Евгеньевич & Работа в связке с задачами Артёма: помощь в реализации компонентов визуализации и интерпретации данных, участие в подготовке демонстрационных кейсов, а также проведение замеров результатов работы системы. \\
\hline
Голод Герман Матвеевич & Реализация модуля SQL-аналитики и запросов к базе данных, включая формирование, выполнение и обработку результатов SQL-запросов. \\
\hline
\end{tabularx}

\section{Ожидаемые результаты}

По итогам выполнения проекта планируется получить следующие результаты:
\begin{enumerate}
    \item рабочую систему «Агентный аналитик данных»;
    \item оркестратор с набором сценариев выполнения;
    \item набор базовых инструментов и подагентов для работы с датасетами;
    \item модуль Data Anonymizer, встроенный в общий конвейер;
    \item экспериментальные данные по времени выполнения и устойчивости сценариев;
    \item оформленные результаты замеров и аналитические выводы;
    \item задел для последующего развития в сторону расширенной мультиоркестрации.
\end{enumerate}

\section{Направления дальнейшего развития}

После завершения первого этапа возможны следующие направления развития:
\begin{itemize}
    \item расширение набора поддерживаемых источников данных;
    \item переход от фиксированных сценариев к более гибкому роутингу;
    \item расширение шаблонного компоновщика промптов;
    \item внедрение полуавтоматического подключения новых инструментов по декларативным схемам;
    \item развитие механизмов песочницы и безопасного выполнения;
    \item переход к более сложной автооркестрации после стабилизации базового ядра.
\end{itemize}


\section{Заключение}

В пояснительной записке зафиксирован подход к реализации проекта «Агентный аналитик данных» как инженерно оправданной и выполнимой системы в рамках проектной практики. В качестве основы выбран оркестратор, так как он обеспечивает управляемость, воспроизводимость и приемлемую сложность реализации. В состав системы включены работа с датасетами, агентно-инструментальный контур аналитики и модуль Data Anonymizer.

\end{document}
